\section{Introduction}\label{sec:introduction}

A central aim of abundance theory is to recover a holomorphic map from
the positivity of an adjoint class. For an ordinary log canonical line,
the expected conclusion is semiampleness of that line itself.
Adding an arbitrary nef line introduces a different phenomenon: its
flat part need not be torsion. The natural conclusion is then numerical
semiampleness, namely the existence of a semiample line with the same
real Chern class.

We prove this statement for smooth compact K\"ahler manifolds with a
klt boundary. A \emph{rational line bundle} is an element of
\(\Pic(X)\otimes_{\Z}\Q\). It is nef when its first Chern class belongs
to the closure of the K\"ahler cone, and pseudo-effective when that
class contains a closed positive \((1,1)\)-current. It is semiample
when a positive integral multiple is a holomorphic line bundle
generated everywhere by global sections.

\begin{theorem}\label{thm:main}
Let \(X\) be a smooth connected compact K\"ahler manifold, let
\(B\geq0\) be a rational simple normal crossing divisor whose
coefficients are strictly less than one, and let
\(M\in\Pic(X)\otimes_{\Z}\Q\).
Assume that \(K_X+B\) is pseudo-effective, \(M\) is nef, and
\(D=K_X+B+M\) is nef. Then there is a semiample rational line bundle
\(L\) on \(X\) such that
\[
 c_1(L)=c_1(D)\quad\text{in }H^{1,1}_{\mathrm{BC}}(X,\R).
\]
\end{theorem}

The numerical formulation is essential even on an elliptic curve:
a nontorsion degree-zero line \(M\) is nef but has no section in any
positive power, whereas its Chern class is represented by the trivial
line. Theorem~\ref{thm:main} allows exactly this change of flat part.
It keeps the nef line \(M\) on the original manifold; it is not a
statement about an arbitrary nef b-divisor on a higher model.

\subsection{Context and the inputs to the proof}

Lazi\'c and Peternell formulated generalized abundance for projective
klt pairs as precisely this numerical semiampleness problem
\cite[Generalised Abundance Conjecture]{LP1}. Their work relates it to ordinary abundance
and to positivity on varieties with numerically trivial canonical
class, using nef reduction and the positive part of the divisorial
Zariski decomposition \cite{LP1,LP2}. They proved the surface case and
the threefold case in which the ordinary adjoint has positive numerical
dimension \cite[Corollaries~C and~D]{LP1}.
The present proof uses the projective positive-part theorem of
\cite[Proposition~8.1]{P} as an input. Our task is to pass from that
projective statement to the analytic setting, where the nef summand
need not be represented by a divisor and a manifold may have no
nonconstant meromorphic functions.

The compact K\"ahler minimal model program has developed alongside
these abundance questions. H\"oring and Peternell constructed minimal
models for compact K\"ahler threefolds \cite{HoeringPeternell2016}.
The threefold abundance theorems
\cite{CHP2016,CHP2023} and the log abundance theorem of Das and Ou
\cite{DasOu2024,DasOu2026} establish important cases of the ordinary problem.
Hacon and Xie's analytic cone theorem \cite{HX} supplies the rational
curve length bound. The restricted K\"ahler model constructions and
ordinary scaling in \cite[Section~3]{K} supply the detected projective
steps used here. For ordinary adjoints in arbitrary dimension, we use
that paper's divisorial decomposition theorem, with its precise
logarithmic Iitaka hypothesis supplied by \cite[Corollary~6.2]{LI}.

These are substantial inputs. Section~\ref{sec:preliminaries} states
their precise forms, retaining the distinction between actual
rational line identities and equalities of real classes.
The proof below supplies the further arguments needed for the nef
summand. It uses Boucksom's analytic divisorial decomposition
\cite{Boucksom2004}, the compactness theory of cycle spaces
\cite{Barlet,LiebermanCycles,Fujiki}, and the geometry of compact K\"ahler
spaces of algebraic dimension zero
\cite{COP,BBYv2,Ou}. The adjoint construction on a projective base
follows the discriminant and moduli-line strategy of Ambro
\cite[Theorem~0.2]{Ambro}, using the period results of \cite{LI}
to accommodate a real polarization.

\subsection{The stronger statement and the proof}

For a pseudo-effective real \((1,1)\)-class \(\alpha\), denote by
\(N(\alpha)\) its divisorial negative part: the effective real divisor
whose coefficient along each prime is the least generic multiplicity
forced in positive representatives, with arbitrarily small K\"ahler
perturbations. A nef class has zero negative part.
We prove a stronger statement in which the adjoint sum need not be nef.

The algebraic dimension \(a(T)\) is the transcendence degree of the
field of meromorphic functions on \(T\). Its algebraic reduction,
after modification, is a morphism \(h:T\to W\) to a smooth projective
variety of dimension \(a(T)\) whose meromorphic functions account for
all those on \(T\).

\begin{theorem}\label{thm:decomposition}
Let \(T\) be a smooth connected compact K\"ahler manifold, \(B\geq0\)
a rational simple normal crossing divisor with coefficients less than
one, and \(M\) a nef rational line bundle. Suppose
\(J=K_T+B\) is pseudo-effective. There are a smooth compact K\"ahler
modification \(\mu:U\to T\), a semiample rational line \(P\) on \(U\),
and an effective rational divisor \(R\) on \(U\), such that
\[
 \mu^*(J+M)\num P+R,\qquad
 R=N\bigl(c_1(\mu^*(J+M))\bigr).
\]
Here \(\num\) denotes equality of real Chern classes.
Moreover, if \(a(T)=0\), then \(c_1(M)=0\).
\end{theorem}

We prove both assertions of Theorem~\ref{thm:decomposition}
simultaneously by induction on dimension. Projective manifolds are
covered by \cite{P}. For \(a(T)=0\), the ordinary decomposition and
the decomposition theorem for K\"ahler klt pairs of Calabi--Yau type
\cite[Corollary~1.4]{BBYv2} reduce the problem to tori and simple
spaces carrying a generically symplectic form. The simple case requires
a meromorphic nonvanishing statement for \(K_T+M\).
Section~\ref{sec:zero} develops the required extension of the
two-diagonal argument in \cite[Section~6]{K}. The new geometric input
excludes a family of correspondences sweeping \(T\times T\): such a
family would produce a nonzero holomorphic one-form on a fixed finite
cover.

When \(0<a(T)<\dim T\), the first objective is to descend \(M\)
numerically to \(W\). Induction decomposes the adjoint on a very general
fiber, but the associated flat twist on that fiber is not known to
extend. Section~\ref{sec:cycles} resolves this difficulty using
countably many global maps constructed from relative cycles.
Each map permits one fixed global choice of flat twist. Coherent
base change then turns positive fiberwise Iitaka dimension into a
meromorphic function contradicting the algebraic reduction.
It follows that \(M\) has zero class on a very general fiber, and
the current-theoretic descent proved in Section~\ref{sec:descent}
gives a nef rational line on \(W\).

It remains to put the ordinary adjoint on a projective base.
After adding a large multiple of an ample line from \(W\), an ordinary
negative program can be run entirely over \(W\).
Its semiample fibration yields
\[
 f:Y\to Z,\qquad K_Y+\Delta\qsim f^*H,
\]
where \(Z\) is projective and \((Y,\Delta)\) is effective and klt.
Section~\ref{sec:adjoint} proves that \(H\) itself is an ordinary
klt adjoint on \(Z\). The key infinitesimal observation identifies
the rank of the extreme Hodge-line period map with the full period
rank. The rational period quotient then supplies the positivity
needed to choose an effective klt boundary on \(Z\).

The projective theorem now applies to the descended adjoint and nef
summand. Section~\ref{sec:completion} identifies its pulled-back
exceptional divisor with the entire analytic negative part using
the mixed Hodge--Riemann relations. This finishes the induction.
For a nef original sum, the negative part vanishes; invariance of
\(\Pic^0\) under smooth modifications descends the semiample
representative and proves Theorem~\ref{thm:main}.

The descent of rational nef classes, the countable construction of
global twists, and the adjoint formula for a K\"ahler total space
are stated separately because each can be used beyond this induction.
